% !TeX root = ../../main.tex
\section{Single-trajectory statistical guarantees}\label{app:stat}

We prove the statistical guarantees of Theorem~\ref{thm:e2e}, using least
squares for CE and a fixed-stage Stein bound for NPG.
First failure combines adaptive stages, and the final count includes
evaluation, baseline, and burn-in transitions.

Put $\kappa=\sqrt{\chi}$. Before sampling or fixing numbered sample
constants, use Lemma~\ref{app:lem:cert} with noise $W+\nu^2BB^\top$
to fix $\bar\sigma^2\ge1$ bounding state and innovation covariances
at every pulse phase. Also fix the LS threshold $L_{\mathrm{LS},0}$
and coefficient $L_{\mathrm{LS}}\ge1$, an absolute logarithm constant
$L_{\log}\ge e$, and a system constant $L_{\mathrm{stat}}\ge1$ dominating
the Gram and Stein bounds.

\subsection{Least squares}\label{app:ls}

Apply $u_t=-K_0x_t+\xi_t$ with independent
$\xi_t\sim N(0,\nu^2I_m)$, independent of process noise.
The counted burn-in is the smallest integer $b_{\mathrm{CE}}\ge1$
with $\kappa\rho^{b_{\mathrm{CE}}}\|x_0\|\le1$, hence
$b_{\mathrm{CE}}\lesssim\log(2+\|x_0\|)$.

\begin{lemma}[Least-squares accuracy]\label{app:lem:ls}
For $0<\delta<1$, collect the preceding experiment with
$N\ge L_{\mathrm{LS},0}(d+\log(1/\delta))$ regression transitions.
With probability at least $1-\delta$, its least-squares fit obeys
\[
 \max\{\|\widehat A-A\|,\|\widehat B-B\|\}^2
 \le L_{\mathrm{LS}}\frac{d+\log(1/\delta)}N .
\]
\end{lemma}
\begin{proof}
\emph{(a) Gaussian design.}
Fix the observed $x_0$, include burn-in innovations, and define
on the regression interval
\[
 z_t=\begin{bmatrix}x_t\\u_t\end{bmatrix},\qquad
 V_N=\sum_tz_tz_t^\top,\qquad
 S_N=\sum_tw_tz_t^\top .
\]
Since $b_{\mathrm{CE}}\ge1$, contemporaneous $x_t,\xi_t$ are
independent and $\operatorname{Cov}(x_t)\succeq W$. The shear
\[
 z_t=\begin{bmatrix}I&0\\-K_0&I\end{bmatrix}
       \begin{bmatrix}x_t\\\xi_t\end{bmatrix}
\]
and its inverse have bounded norms, giving uniformly positive
design covariance bounds. Burn-in bounds the deterministic mean.
The convolution bound
$\sum_{j\ge0}\|F_{K_0}^j\|\le\kappa/(1-\rho)$ bounds the complete
Gaussian stack covariance by $\lesssim1$, including burn-in,
exploration, and shear.

Thus $(v^\top z_t)_t$ has covariance norm $\lesssim1$ and second
moment $\lesssim N$ for unit $v\in\RR^d$.
Diagonalization and Gaussian moment-generating functions give squared
norm deviations $\lesssim\sqrt{Nu}+u$ with failure $O(e^{-u})$.
A $1/4$-net of size $\exp(O(d))$, with
$u\gtrsim d+\log(1/\delta)$, yields
\[
 \left\|N^{-1}(V_N-\EE V_N)\right\|
 \lesssim\sqrt{\frac{d+\log(1/\delta)}N}
              +\frac{d+\log(1/\delta)}N .
\]
The sample threshold gives
$\lambda_{\min}(V_N)\gtrsim N$ and $\|V_N\|\lesssim N$.

\emph{(b) Martingale noise.}
Since $z_t$ precedes $w_t$, apply the self-normalized inequality
\citep{abbasi2011improved} to $S_N^\top v$ for fixed $v\in\RR^n$
with regularizer $\lambda NI$, where $\lambda>0$ is fixed.
On the Gram event, the log determinant is $\lesssim d$ and
$\|V_N^{-1/2}(V_N+\lambda NI)V_N^{-1/2}\|\lesssim1$.
A sphere net adds $O(n)$ to the logarithm, giving
$\|S_NV_N^{-1/2}\|^2\lesssim d+\log(1/\delta)$.
Intersect the two events within failure $\delta$ and use
$[\widehat A-A,\widehat B-B]=S_NV_N^{-1}$.
The desired claim then follows.
\end{proof}

Fix
\[
 C_3=\max\{\Lambda_{\mathrm{CE}},\,2C_8/c_2\}.
\]
Fix $C_2$ ensuring its sample condition implies
$N\ge L_{\mathrm{LS},0}(d+\log(1/\delta))$,
$L_{\mathrm{LS}}(d+\log(1/\delta))/N\le c_1^2$, and
$\Lambda_{\mathrm{CE}}rL_{\mathrm{LS}}(d+\log(1/\delta))/N\le\eps$.
Fix $C_4$ satisfying, for $n\ge1$ and $0<\eps<1$,
\[
 1+\frac2{c_2}\log(2+2nw_+p_0/\eps)
 \le C_4\log(C_4n/\eps).
\]

\begin{proof}[Proof of Theorem~\ref{thm:e2e}(a)]
\emph{(a) Accuracy.}
By Lemma~\ref{app:lem:ls} and the choice of $C_2$, with
probability at least $1-\delta$,
\[
 a=\sqrt{L_{\mathrm{LS}}(d+\log(1/\delta))/N}
 \le c_1\le e_{\mathrm{loc}}/3 .
\]

\emph{(b) Transfer.}
Since $\beta_r\le c_1$, Proposition~\ref{app:prop:ce}
gives stability, matrix and rank conclusions, and
\[
 g(\widehat K)\le\Lambda_{\mathrm{CE}}(ra^2+\tau_r^2)
 \le\eps+C_3\tau_r^2.
\]
The total transition count is $N+b_{\mathrm{CE}}$.
The desired claim then follows.
\end{proof}

\subsection{A fixed Stein stage}\label{app:stein}

Condition on a stage-entry history, fixing its state $x_{\mathrm{in}}$
and stabilizing gain $K$ with $P=P_K\preceq2P_{K_0}$.
Write $F=F_K$ and choose an integer $H\ge1$ satisfying
\begin{equation}\label{app:eq:horizon}
 \kappa\rho^H\sqrt{\bar\sigma^2/w_-}\le\tfrac12.
\end{equation}
Every block, including burn-in, begins with an independent pulse
$\xi_j\sim N(0,\nu^2I_m)$ and continues with $H-1$ nominal
transitions. Choose the smallest nonnegative burn-in count with
\begin{equation}\label{app:eq:burn}
 \kappa^2\rho^{2b_{\mathrm{burn}}H}
       (\|x_{\mathrm{in}}\|^2+n\bar\sigma^2)\le1.
\end{equation}
After burn-in, set $t_j=jH$, use baseline block $0$ and evaluation
blocks $1,\ldots,N$, and define
\[
 \begin{gathered}
 X_j=x_{t_j},\qquad
 \mathcal G_j=\sum_{s=1}^H x_{t_j+s}^\top S_Kx_{t_j+s},\qquad
 \Delta_j=\mathcal G_j-\mathcal G_{j-1},\\
 Y_j=-\frac{\xi_j\Delta_j}{2\nu^2},\qquad
 T_j=\frac{(\xi_j\xi_j^\top-\nu^2I_m)\Delta_j}{2\nu^4},\qquad
 \widehat G=N^{-1}\sum_{j=1}^NT_j .
 \end{gathered}
\]
Set
\[
 \widehat Z=\Bigl(\sum_{j=1}^NY_jX_j^\top\Bigr)
             \Bigl(\sum_{j=1}^NX_jX_j^\top\Bigr)^\dagger,
\]
which is \eqref{eq:stein}, with the minimum-norm convention for a
singular Gram matrix. The return excludes immediate pulse
cost and ends with nominal cost at the endpoint before the next pulse.
For the pre-pulse history $\mathcal F_j$, both
$X_j,\mathcal G_{j-1}$ are measurable.
The stage counts $(b_{\mathrm{burn}}+N+1)H$ transitions.

\begin{lemma}[Relative state Gram matrix]\label{app:lem:gram}
Suppose $K$ is stabilizing with $P_K\preceq2P_{K_0}$ and use the
stage protocol with \eqref{app:eq:horizon}--\eqref{app:eq:burn}.
If $0<\delta<1$ and $N\ge L_{\mathrm{stat}}(n+\log(1/\delta))$, then,
given the stage-entry history, with probability at least $1-\delta$,
\[
 \tfrac12N\Sigma_H\preceq\sum_{j=1}^NX_jX_j^\top\preceq2N\Sigma_H,
 \qquad w_-I\preceq\Sigma_H\preceq\bar\sigma^2I .
\]
Here $\Sigma_H$ is the stationary block-start covariance.
\end{lemma}
\begin{proof}
\emph{(a) Normalization.}
The boundary process $X_{j+1}=F^HX_j+\omega_j$ has independent
Gaussian innovations with covariance $\succeq W$ and stationary
state covariance $\Sigma_H\preceq\bar\sigma^2I$.
For $\overline X_j=\Sigma_H^{-1/2}X_j$ and
$T_H=\Sigma_H^{-1/2}F^H\Sigma_H^{1/2}$, the horizon rule gives
$\|T_H\|\le1/2$, hence normalized innovation covariance
$I-T_HT_H^\top\succeq3I/4$.
Summing lag bounds $2^{-a}$ bounds the stationary stack covariance
by $O(1)$. This dominates the centered stack from a fixed state
under the same burn-in protocol.
The baseline block gives each evaluation covariance $\succeq3I/4$,
while \eqref{app:eq:burn} gives
$\sum_j\|\EE\overline X_j\|^2\lesssim1$.

\emph{(b) Concentration.}
Gaussian diagonalization and linear-form tails give, for
$g\sim N(0,\mathsf V)$ and fixed $\mu$,
\[
 \bigl|\|g+\mu\|^2-\tr\mathsf V-\|\mu\|^2\bigr|
 \lesssim \|\mathsf V\|_F\sqrt{u}+\|\mathsf V\|u
                       +\sqrt{\mu^\top\mathsf V\mu\,u}
\]
with failure $O(e^{-u})$ \citep{laurent2000adaptive,rudelson2013hanson}.

A fixed-resolution net applied to $(v^\top\overline X_j)_j$ gives
normalized Gram deviation $\lesssim\sqrt{u/N}+u/N+\sqrt{u}/N$
for $u\gtrsim n+\log(1/\delta)$.
By part (a), its expectation lies between $3NI/4$ and
$NI+\sum_j(\EE\overline X_j)(\EE\overline X_j)^\top=NI+O(1)$.
The fixed $L_{\mathrm{stat}}$ and undoing normalization give the bounds.
The desired claim then follows.
\end{proof}

\begin{lemma}[Stein accuracy]\label{app:lem:stein}
Suppose $K$ is stabilizing with $P_K\preceq2P_{K_0}$ and use the
stage protocol with \eqref{app:eq:horizon}--\eqref{app:eq:burn}.
For $0<\delta<1$, let $\ell=\log(L_{\log}Nd/\delta)$.
If $N\ge L_{\mathrm{stat}}d\ell^3$, then, given the stage-entry history,
with probability at least $1-\delta$,
\[
 \|[\widehat G-G_K,\widehat Z-Z_K]\|^2
 \le L_{\mathrm{stat}}\left(\frac{Hnd\ell^3}{N}+\rho^{4H}\right).
\]
The relative Gram conclusion of Lemma~\ref{app:lem:gram} holds
on the same event.
\end{lemma}
\begin{proof}
\emph{(a) Targets and finite-horizon bias.}
Put $P_H=\sum_{a=0}^{H-1}(F^a)^\top S_KF^a$,
$G_H=B^\top P_HB$, and $Z_H=-B^\top P_HF$.
The within-block state expansion gives
\[
 \EE[\mathcal G_j\given\mathcal F_j,\xi_j]
 =X_j^\top F^\top P_HFX_j+2\xi_j^\top B^\top P_HFX_j
             +\xi_j^\top G_H\xi_j+\gamma_H ,
\]
where $\gamma_H$ is independent of $X_j,\xi_j$.
Since $\mathcal G_{j-1}$ is $\mathcal F_j$-measurable and
$\EE[\xi_j\given\mathcal F_j]=0$,
$\EE[\xi_j\xi_j^\top-\nu^2I_m\given\mathcal F_j]=0$, replacing
$\mathcal G_j$ by $\Delta_j$ leaves both conditional targets unchanged. Vanishing odd
moments and
$\EE[(\xi\xi^\top-\nu^2I)(\xi^\top G_H\xi)]=2\nu^4G_H$ give
\begin{equation}\label{app:eq:targets}
 \EE[Y_j\given\mathcal F_j]=Z_HX_j,\qquad
 \EE[T_j\given\mathcal F_j]=G_H.
\end{equation}
Since $P-P_H=(F^H)^\top PF^H$,
\begin{equation}\label{app:eq:bias}
 \|G_H-G_K\|+\|Z_H-Z_K\|\lesssim\rho^{2H}.
\end{equation}

\emph{(b) Return second moments.}
Let $V_a=(F^a)^\top PF^a$. Summation by parts, using
$V_{a+1}\preceq\rho^2V_a$, gives
\[
 \sum_{a\ge0}\rho^{-a}(F^a)^\top S_KF^a\preceq(1+\rho)P .
\]
Weighted Cauchy--Schwarz bounds
$\|S_K^{1/2}\sum_{a\ge0}z^aF^a\|^2\le p(1+\rho)/(1-\rho)$
for $|z|=1$. Parseval bounds the innovation-to-cost-weighted-state
map. Every stationary phase covariance is dominated by a system multiple
of $\sum_aF^a(F^a)^\top$, whose cost-weighted trace is
$\lesssim\tr P$.

Under the fixed stage-entry history, stack the cost-weighted states of
the two returns as
\[
 y_j=\operatorname{col}\bigl(S_K^{1/2}x_{t_{j-1}+s}\bigr)_{s=1}^{2H},
 \qquad \mu_y=\EE y_j,\qquad \mathsf V_y=\operatorname{Cov}(y_j),
\]
so that $\Delta_j=y_j^\top\mathsf Jy_j$ with
$\mathsf J=\operatorname{diag}(-I_{Hn},I_{Hn})$. Then
$\|\mathsf V_y\|\lesssim1$, $\tr\mathsf V_y\lesssim Hn$, and
$\|\mu_y\|^2\lesssim1$ by Lyapunov energy and
\eqref{app:eq:burn}. The stationary returns have equal means.
Under the same pulsed burn-in, the first-boundary covariance deficit
is $F^{b_{\mathrm{burn}}H}\Sigma_H(F^{b_{\mathrm{burn}}H})^\top$,
giving
\[
 |\EE\Delta_j|\lesssim\kappa^2\rho^{2b_{\mathrm{burn}}H}
 (\|x_{\mathrm{in}}\|^2+n\bar\sigma^2)\lesssim1 .
\]

The Gaussian quadratic-form identity gives
\[
 \EE\Delta_j^2
 =2\tr(\mathsf J\mathsf V_y\mathsf J\mathsf V_y)
  +4\mu_y^\top\mathsf J\mathsf V_y\mathsf J\mu_y+(\EE\Delta_j)^2
 \lesssim Hn .
\]
Write $V_H=Hn$. The same calculation gives
$\operatorname{Var}(\mathcal G_j)\lesssim V_H$.

\emph{(c) Predictable moments.}
Write
$d_j=\EE[\Delta_j\given\mathcal F_j]$,
$\varsigma_j^2=\operatorname{Var}(\mathcal G_j\given\mathcal F_j)$, and
$m_j=d_j^2+\varsigma_j^2=\EE[\Delta_j^2\given\mathcal F_j]$,
where the last equality uses the $\mathcal F_j$-measurability of
$\mathcal G_{j-1}$.
Here $d_j$ has Gaussian degree at most two, and
$\varsigma_j^2$ is a nonnegative constant plus a nonnegative quadratic
form in $X_j$. Jensen's inequality and the law of total variance give
$\EE d_j^2,\EE\varsigma_j^2\lesssim V_H$.

For a degree-$k$ Gaussian polynomial $f$,
$\|f\|_{L^a}\le(a-1)^{k/2}\|f\|_{L^2}$ for $a\ge2$
\citep{nelson1973free}.
Applied to $d_j$ and the sum-of-squares representation of
$\varsigma_j^2$, this gives $\|m_j\|_{L^a}\lesssim a^2V_H$
by Minkowski. Markov at moment of order $\ell$ and a union bound yield
\begin{equation}\label{app:eq:mh}
 \max_{j\le N}m_j\lesssim\ell^2V_H
\end{equation}
with the allocated failure probability.

Let $\widetilde Y_j=Y_j-Z_HX_j$ and $D_j^G=T_j-G_H$.
Conditional H\"older's inequality and degree-two hypercontractivity
of $\Delta_j$ give, for unit $u,v$,
\[
 \begin{split}
 u^\top\EE[\widetilde Y_j\widetilde Y_j^\top\given\mathcal F_j]u
       &\lesssim m_j,\\
 \EE[(u^\top D_j^Gv)^2\given\mathcal F_j]&\lesssim m_j,\qquad
 \|\EE[(D_j^G)^2\given\mathcal F_j]\|\lesssim m\,m_j .
 \end{split}
\]
The last bound sums over one action index.

\emph{(d) Predictable quadratic variation.}
Set $\overline M_j=\widetilde Y_j\overline X_j^\top$ using the
normalization of Lemma~\ref{app:lem:gram}.
These are martingale differences from $\mathcal F_j$ to
$\mathcal F_{j+1}$ by \eqref{app:eq:targets}.
On the moment and upper Gram events, the variations satisfy
\[
 v_Z\lesssim N\ell^2V_H\max\{n,m\},\qquad
 v_G\lesssim Nm\ell^2V_H
\]
for $\sum_j\overline M_j$ and $\sum_jD_j^G$, respectively.
For the first sequence, the conditional matrix products are
$\|\overline X_j\|^2\EE[\widetilde Y_j\widetilde Y_j^\top\given\mathcal F_j]$
and
$\tr\EE[\widetilde Y_j\widetilde Y_j^\top\given\mathcal F_j]\,
\overline X_j\overline X_j^\top$.
The normalized Gram trace and norm bound their sums.
Increment entries have Gaussian degree at most four and
$L^2$ norms $\lesssim\sqrt{V_H}$ by H\"older.
Hypercontractivity and $\|M\|\le\|M\|_F$ give, for $a\ge2$,
\begin{equation}\label{app:eq:tails}
 \bigl\|\|\overline M_j\|\bigr\|_{L^a}
      \lesssim\sqrt{mnV_H}\,a^2,\qquad
 \bigl\|\|D_j^G\|\bigr\|_{L^a}
      \lesssim m\sqrt{V_H}\,a^2 .
\end{equation}

\emph{(e) Bounded increments and compensators.}
For $(U_j,v)=(\overline M_j,v_Z)$ or $(D_j^G,v_G)$, let $A_U$ be the
deterministic coefficient in \eqref{app:eq:tails}, so that
$\|\|U_j\|\|_{L^a}\le A_Ua^2$ for all $a\ge2$. Put
$L_{\mathrm{cut}}=e^2A_U(4\ell)^2$ and
\[
 \widetilde U_j
 =U_j\mathbf1_{\{\|U_j\|\le L_{\mathrm{cut}}\}}
 -\EE[U_j\mathbf1_{\{\|U_j\|\le L_{\mathrm{cut}}\}}\given\mathcal F_j].
\]
These martingale differences have norm $\le2L_{\mathrm{cut}}$
and conditional second moments dominated by those of $U_j$.
The $4\ell$th moment in \eqref{app:eq:tails} gives
\[
 \Pr(\|U_j\|>L_{\mathrm{cut}})\le e^{-4\ell},\qquad
 \EE[\|U_j\|\mathbf1_{\{\|U_j\|>L_{\mathrm{cut}}\}}]
           \le L_{\mathrm{cut}}e^{-3\ell}.
\]
A union bound and Markov control observed truncations and conditional
tail compensators, giving
$\|\sum_j(U_j-\widetilde U_j)\|\le L_{\mathrm{cut}}$
outside probability $N(e^{-4\ell}+e^{-3\ell})$.

Matrix Freedman \citep{tropp2011freedman}, with self-adjoint
dilation for rectangular increments, gives
\[
 \left\|\sum_{j=1}^NU_j\right\|
 \lesssim\sqrt{v\ell}+L_{\mathrm{cut}}\ell
\]
on the variation event, outside the allocated exceptional event.

\emph{(f) Reconstruction.}
Put $\overline\Sigma_N=N^{-1}\sum_j\overline X_j\overline X_j^\top$.
Regression gives
\[
 \widehat Z-Z_H
 =\left(N^{-1}\sum_j\overline M_j\right)
    \overline\Sigma_N^{-1}\Sigma_H^{-1/2}.
\]
The last factors are bounded by $2$ and $1/\sqrt{w_-}$
by Lemma~\ref{app:lem:gram}. Substituting variations and cutoffs gives
\[
 \begin{split}
 \|\widehat Z-Z_H\|
 &\lesssim\sqrt{\frac{V_H\max\{n,m\}}N}\,\ell\sqrt{\ell}
           +\frac{\sqrt{mnV_H}}N\ell^3,\\
 \|\widehat G-G_H\|
 &\lesssim\sqrt{V_H}
          \left\{\sqrt{\frac mN}\,\ell^{3/2}
                        +\frac mN\ell^3\right\}.
 \end{split}
\]
Since $N\gtrsim d\ell^3$, the second terms are absorbed.
Use $\|[E_G,E_Z]\|^2\le\|E_G\|^2+\|E_Z\|^2$,
$V_H=Hn$, and \eqref{app:eq:bias}, with total failure $\delta$.
The desired claim then follows.
\end{proof}

\subsection{Schedule for NPG}\label{app:schedule}

With $c_1,c_2,C_8$ from Appendix~\ref{app:transfer}, set
\[
 \bar g_t=nw_+p_0(1-c_2/2)^t,\qquad
 a_t^2=\min\left\{c_1^2,\frac{c_2\bar g_t}{4C_8r}\right\},
 \qquad T=\min\{t\ge0\given\bar g_t\le\eps\}.
\]
When $T=0$, return $K_0$.
For $T>0$, choose the smallest nondecreasing integer horizons $H_t$
satisfying \eqref{app:eq:horizon}, $\kappa^2\rho^{H_t}\le1$, and
$L_{\mathrm{stat}}\rho^{4H_t}\le a_t^2/2$.

Put $L_{\mathrm{batch}}=2L_{\mathrm{stat}}$ and $L_*=2L_{\log}$.
Before sampling, choose by deterministic doubling a common $N_*$
such that, with $\ell_*=\log(L_*N_*Td/\delta)$,
\[
 N_*\ge
 \left\lceil L_{\mathrm{batch}}
 \left(d\ell_*^3+\frac{H_{T-1}nd\ell_*^3}{a_{T-1}^2}\right)\right\rceil.
\]
Use
\begin{equation}\label{app:eq:batches}
 N_t=\left\lceil L_{\mathrm{batch}}
      \left(d\ell_*^3+\frac{H_tnd\ell_*^3}{a_t^2}\right)\right\rceil
 \le N_* .
\end{equation}
For $U_*=d+H_{T-1}nd/a_{T-1}^2$, doubling terminates since
$N/\log^3N\to\infty$, with
\[
 N_*\lesssim U_*\log^3(2U_*Td/\delta).
\]
Thus $\ell_*$ grows logarithmically. For $\delta_t=\delta/(2T)$, the
stage logarithm satisfies
\[
 \ell_t=\log(L_{\log}N_td/\delta_t)\le\ell_*,
\]
and the batch and horizon rules give
\[
 N_t\ge L_{\mathrm{stat}}d\ell_t^3,\qquad
 L_{\mathrm{stat}}\Bigl(\frac{H_tnd\ell_t^3}{N_t}+\rho^{4H_t}\Bigr)
 \le a_t^2.
\]
Lemma~\ref{app:lem:stein} therefore gives radius $a_t$ with conditional
failure at most $\delta_t$ on certified histories.

\subsection{Adaptive stages and total burn-in}\label{app:stages}

For the pre-stage history $\mathcal H_t$, let $E_t$ mean
$\|\widehat Z_t-Z_{K_t}\|\le a_t$ on certified histories, and be
false otherwise. Then $\mathcal A_t=\bigcap_{s<t}E_s\in\mathcal H_t$.
Induction using Proposition~\ref{app:prop:npg} certifies $K_t$ on
$\mathcal A_t$, where Lemma~\ref{app:lem:stein} gives
$\Pr(E_t^c\given\mathcal H_t)\le\delta_t$.
The disjoint first-failure events satisfy
\begin{equation}\label{app:eq:first-failure}
 \Pr(\mathcal A_t\cap E_t^c)
 =\EE[\mathbf1_{\mathcal A_t}\Pr(E_t^c\given\mathcal H_t)]
 \le\delta_t,\qquad
 \Pr\left(\bigcap_{t<T}E_t\right)\ge1-\sum_{t<T}\delta_t.
\end{equation}


\begin{lemma}[Total burn-in]\label{app:lem:burn}
Let $T,N_t,H_t$ be positive integers, with deterministic $N_t,H_t$,
and use \eqref{app:eq:burn} at every stage. Suppose
$\kappa^2\rho^{H_t}\le1$ for every $t$.
With probability at least $1-\delta$, on every prefix of the realized
trajectory whose gains satisfy the stage certificates,
\[
 \sum_{t<T}b_{{\mathrm{burn}},t}H_t-\sum_{t<T}H_t
 \lesssim T\log\frac{2(1+\|x_0\|)n(M+1)}\delta .
\]
Here $M=\sum_{t<T}(b_{{\mathrm{burn}},t}+N_t+1)H_t$ is the actual
transition count.
\end{lemma}
\begin{proof}
\emph{(a) A pathwise state envelope.}
Full blocks contract by $\kappa\rho^{H_t}\le\rho^{H_t/2}$,
and length-$a$ partial blocks by $\kappa\rho^{a/2}$.
Any transition product has at most two partial endpoints, hence
$\|\Phi(v,u)\|\le\kappa^2\rho^{(v-u)/2}$ on certified prefixes,
including policy changes.

The noise $\omega_s=w_s+B\xi_s$ when pulsed, and $w_s$ otherwise,
has bounded conditional covariance.
Gaussian tails with failure allocation proportional to
$\delta/(s+1)^2$ give, simultaneously with probability $1-\delta$,
\[
 \|\omega_s\|\lesssim\sqrt{n}+
            \sqrt{\log\frac{O(1)(s+1)^2}\delta}
\]
for all $s\ge0$.
Variation of constants and the summable transition bound yield
\[
 \kappa^2(\|x_s\|^2+n\bar\sigma^2)
 \lesssim \|x_0\|^2+n+\log\frac{O(1)(M+1)}\delta,\qquad s\le M.
\]

\emph{(b) The integer rule.}
Minimality and $-\log\rho\ge1/(2\chi)$ give
\[
 b_{{\mathrm{burn}},t}H_t
 \le H_t+\chi\log_+\{\kappa^2(\|x_{{\mathrm{in}},t}\|^2+n\bar\sigma^2)\},
\]
where $\log_+z=\max\{\log z,0\}$.
Insert the state envelope and sum over $t$.
The desired claim then follows.
\end{proof}

\begin{proof}[Proof of Theorem~\ref{thm:e2e}(b)]
For $T=0$, $g(K_0)\le\bar g_0\le\eps$ and $M=0$.
Otherwise proceed as follows.

\emph{(a) Probability and containment.}
Intersect \eqref{app:eq:first-failure} and Lemma~\ref{app:lem:burn},
each at failure $\delta/2$.
With probability $1-\delta$, Proposition~\ref{app:prop:npg}
gives stability and $P_{K_t}\preceq2P_{K_0}$ for every $t\le T$.

\emph{(b) Gap envelope.}
Initially $g(K_0)\le J(K_0)\le nw_+p_0=\bar g_0$.
If $g(K_t)\le\bar g_t+C_3\tau_r^2$, then
\[
 \begin{split}
 g(K_{t+1})
 &\overset{\mathrm{(a)}}{\le}
 (1-c_2)\bar g_t+\frac{c_2}{4}\bar g_t+
            \{(1-c_2)C_3+C_8\}\tau_r^2\\
 &\overset{\mathrm{(b)}}{\le}\bar g_{t+1}+C_3\tau_r^2.
 \end{split}
\]
Here (a) uses the transfer recursion and the definition of $a_t$,
and (b) uses $C_3\ge2C_8/c_2$ and $1-3c_2/4\le1-c_2/2$.
Induction gives the terminal gap. Since $\bar g_{T-1}>\eps$,
\[
 T\le1+2\log(nw_+p_0/\eps)/c_2\le C_4\log(C_4n/\eps),
\]
also valid for $T=0$.

\emph{(c) Transitions.}
Geometric summation gives
\[
 \sum_{t<T}a_t^{-2}
 \lesssim T+r\sum_{t<T}\bar g_t^{-1}
 \lesssim T+r/\eps .
\]
Also $a_t^2\gtrsim\min\{1,\eps/r\}$ for $t<T$, giving
$H_t\lesssim\log(2r/\eps)$.
Writing $H_*=H_{T-1}$, \eqref{app:eq:batches} yields
\[
 \sum_{t<T}H_tN_t
 \lesssim H_*T(d\ell_*^3+1)
              +H_*^2nd\ell_*^3(T+r/\eps).
\]
Adding baseline and burn-in counts gives
$M\lesssim M_{\mathrm{det}}+T\log(2M)$, where
$M_{\mathrm{det}}=\widetilde O_{\sys}(rnd/\eps)$ includes logarithms
of $1+\|x_0\|$ and $1/\delta$.
Elementary inversion gives
$M\lesssim M_{\mathrm{det}}+T\log(2(M_{\mathrm{det}}+T))
=\widetilde O_{\sys}(rnd/\eps)$.
The proof is then completed.
\end{proof}

The dense assertions follow using Proposition~\ref{app:prop:ce}'s
dense bound for CE, and $r=\min\{n,m\}$ with the dense transfer
recursion for NPG. Both tail terms vanish.



